\documentclass[12pt, a4paper, oneside]{book}

\usepackage{graphicx}
\usepackage{amsmath}
\usepackage{amsfonts}
\usepackage{mathtools}
\usepackage{amsthm}
\usepackage{amssymb}
\usepackage{newlfont}
\usepackage{listings}
\usepackage{xcolor}
\usepackage{verbatim}

% TODO QUI È DOVE INIZIA LA DEFINIZIONE DELLE SCRITTURE DI CODICE

\definecolor{codegreen}{rgb}{0,0.6,0}
\definecolor{codegray}{rgb}{0.5,0.5,0.5}
\definecolor{codepurple}{rgb}{0.58,0,0.82}
\definecolor{backcolour}{rgb}{0.95,0.95,0.92}

\lstdefinestyle{mystyle}{
	backgroundcolor=\color{backcolour},   
	commentstyle=\color{codegreen},
	keywordstyle=\color{magenta},
	numberstyle=\tiny\color{codegray},
	stringstyle=\color{codepurple},
	basicstyle=\ttfamily\footnotesize,
	breakatwhitespace=false,         
	breaklines=true,                 
	captionpos=b,                    
	keepspaces=true,                 
	numbers=left,                    
	numbersep=5pt,                  
	showspaces=false,                
	showstringspaces=false,
	showtabs=false,                  
	tabsize=2
}

\lstset{style=mystyle}
\makeatother

\newtheorem{theorem}{Teorema}
\newtheorem{definition}{Definizione}
\newtheorem{eg}{Esempio}
\newcommand*{\blankpage}{
{\newpage \vspace*{5cm}\thispagestyle{empty}\centering \bfseries  \textit{Mamma mia, oggi sembro handicappata - Gloria M.}\par}
\vspace{\fill}}

\textwidth=450pt\oddsidemargin=0pt

\begin{document}
	
	\begin{titlepage}
		\begin{center}
			{{\LARGE{\textsc{Università degli Studi di Verona}}}} \rule[0.1cm]{15.8cm}{0.1mm}
			\rule[0.5cm]{15.8cm}{0.6mm}
			{\small{\bf CORSO DI LAUREA IN INFORMATICA}}
		\end{center}
		\vspace{12.5mm}
		
		\begin{center}
			{\Huge{\bf Elaborazione di segnali e immagini}}
		\end{center}
		\vspace{130mm}
		
		\begin{minipage}[l]{1.8mm}
			\rule[0.1cm]{0.1mm}{2cm}
			\rule[0.1cm]{0.6mm}{2cm}
		\end{minipage}
		\hspace{2mm}
		\begin{minipage}{10cm}	
			{\large{Federico Brutti\\
			\texttt{federico.brutti@studenti.univr.it}}}
		\end{minipage}
	\end{titlepage}
	
	\blankpage
	\tableofcontents
	
	\chapter{Introduzione e signal processing}
	\section{Concetti di base e tipi di segnali}
Il corso di elaborazione dei segnali e immagini provvede a fornire strumenti, sia pratici che teorici, per l'elaborazione di segnali n-dimensionali,
rappresentanti quantità fisiche o strutture visive. Conoscere le caratteristiche dei segnali per elaborare immagini e interpretarne il contenuto è
una competenza che aiuta a interagire e conoscere fenomeni di vario tipo, la quale è tipicamente applicata in vari ambiti: creazione di contenuti
multimediali, ispezioni visive per il controllo qualità, medico, per fare diagnosi, o videosorveglianza. Esistono diversi tipi di segnali, ma tutti
seguono il seguente iter di elaborazione: \begin{itemize}
	\item Ricezione del \textbf{segnale analogico} tramite apposito strumento. Si ottiene una funzione continua $x(t)$.
	\item \textbf{Campionamento}: Divide la funzione in più intervalli di tempo dal	valore discreto $x[t]$.
	\item \textbf{Quantizzazione}: Mappatura di ogni campione ottenuto per arrotondare il valore di ampiezza a uno discreto $x_q[t]$.
\end{itemize}

\noindent Il risultato di questo processo è un segnale digitale, descritto da una sequenza di numeri, sulla quale si può attuare l'iter di elaborazione,
inteso come l'esecuzione di operazioni matematiche su di essa. Dopodiché, sarà necessario riportare il segnale modificato nel dominio del tempo, quindi
in forma analogica, nel rispetto del teorema di Nyquist-Shannon. Oggi ogni computer ha a disposizione vari programmi per la conversione dei segnali, i
cui due tipi principali sono \textbf{unidimensionali} e \textbf{bidimensionali}. I primi descrivono tipicamente quantità fisiche in funzione di una
variabile indipendente, come il tempo. Il campionamento discretizza la variabile indipendente, mentre la quantizzazione si occupa dell'ampiezza del
segnale.\par
I \textbf{segnali 2D}, invece, sono una rappresentazione di una scena o oggetto visuale, la quale è descritta tramite una matrice bidimensionale i cui
elementi, detti \textbf{pixel}, contengono valori numerici relativi a componenti di colore. In questo contesto, campionamento e quantizzazione
concernono lo spazio nelle due dimensioni; a partire da un'immagine analogica, la si passa per il sistema di conversione, il quale la proietterà su una
matrice per poterla campionare\footnote{Fun fact, negli anni '20 si usava fondamentalmente la stessa tecnica e i valori dei "pixel" erano ottenuti
tramite trasmissioni telegrafiche.}. In base al risultato ottenuto, la si discretizzerà determinando i valori relativi al colore contenuti dentro ai
pixel. In particolare, per le immagini digitali definiamo alcune misure legate al campionamento: \begin{itemize}
	\item \textbf{Dimensione}: Numero di pixel disponibili all'interno della matrice, definito come il prodotto tra pixel orizzontali e verticali: \[
		\text{pixel}_H\times \text{pixel}_V
	\]
	\item \textbf{Risoluzione}: Numero di punti generati all'interno di un'area	pre-definita, come per esempio un $cm^2$. Si tratta di un valore legato
	al supporto che genera l'immagine. Più sono, più l'immagine sarà risoluta e più bassa sarà la probabilità di perdere dati.
\end{itemize}

\noindent Come già menzionato, i valori all'interno dei pixel della matrice riguardano le componenti del colore, le quali sono descritte dalla \textbf{
profondità del colore}, legata alla quantizzazione, determina il numero di bit allocati per fare variazioni di colore\footnote{tipicamente sono $8b$ per
canale, quindi un intervallo di $[0, 255]$. Non necessariamente ogni pixel deve avere la stessa profondità}, e dallo \textbf{spazio colore}, una scala
la cui somma dei tre elementi determina il colore effettivo. In particolare, esistono più spazi: \begin{itemize}
	\item \textbf{RGB}: Red Green Blue è la più comune e fa uso della somma dei valori degli omonimi tre colori primari.
	\item \textbf{HSL}: Hue Saturation Lightness determina il colore in base alla tinta (colore effettivo), saturazione (la sua intensità rispetto al
	grigio) e luminosità.
	\item \textbf{CMY}: Cyan, Magenta, Yellow, partendo dal bianco, sottrae i tre valori per creare le varie tinte.
\end{itemize}

\noindent Avere più spazi colori a disposizione è molto importante, poiché consente di effettuare operazioni su un range di caratteristiche più ampio.
In particolar modo si fa uso degli \textbf{istogrammi delle immagini}, i quali agiscono come rappresentazione grafica della distribuzione delle tonalità
dell'immagine, consentendo di giudicare i valori tonali in una grafica più contenuta. Lo si può inoltre fornire ai computer per effettuare elaborazioni
specifiche su una determinata proprietà dell'immagine in questione. Generalmente, si ottengono immagini visivamente accettabili quando l'istogramma è
piatto, poiché avendo più valori contrastanti è più semplice notare le eventuali differenze.\par
Supponiamo ora di avere quattro elementi di colori diversi su uno sfondo omogeneo. Ci è possibile isolarli facilmente grazie all'istogramma, perché avrà
il picco più alto data la prevalenza del colore di sfondo. Rimuovendolo, si isoleranno gli elementi. Altri tipi di segnali affrontati in questo corso
saranno: \begin{itemize}
	\item \textbf{Immagini dinamiche}: Video, composti da un insieme discreto di immagini, campionate in un determinato intervallo di tempo (FPS).
	\item \textbf{Immagini volumetriche}: Volume di un'immagine ottenuto estendendo	la matrice di campionamento in tre dimensioni.
	\item \textbf{3D Mesh}: Visualizzazione della superficie tridimensionale di un modello il cui involucro è mappabile su un'immagine 2D tramite mappa
	UV.
\end{itemize}

%

\section{Modalità di immagini}
Oltre ai tipi di immagine discussi esisono anche varie \textbf{modalità} per la loro acquisizione in base allo strumento o al sensore usato. In
particolare abbiamo: \begin{itemize}
	\item \textbf{Immagini di profondità}: Chiamate anche mappe range, enfatizzano la geometria dell'immagine indicando la distanza dell'oggetto
	ritratto in base al valore dei pixel della matrice. Il valore di profondità è legato al punto di vista dello strumento usato per l'acquisizione, ed
	è definito come la distanza tra punto di vista della camera e il primo punto dell'oggetto incontrato lungo il raggio di proiezione. Tipicamente sono
	usate per la conversione di immagini tridimensionali in modelli 3D.
	\item \textbf{Immagini stereo}: Chiamate anche mappe di disparità, sono il prodotto di due immagini acquisite da due punti di vista diversi,
	tipicamente usate quando si vogliono acquisire informazioni geometriche tridimensionali di una scena intera.\par
	Osservando uno stesso oggetto da due punti di vista diversi si proiettano due rette su di esso, consentendo di disambiguare la posizione esatta del
	punto sull'asse di proiezione. Sfruttano inoltre il concetto di \textbf{parallasse}, che mostra come la distanza varia la velocità di spostamento
	degli oggetti all'interno della scena; in particolare, più vicini sono, più velocemente si sposteranno.
	\item \textbf{Immagini ultrasuoni}: Acquisite tramite elaborazione dei segnali ultrasuoni, codificano la profondità di una scena in base al tempo di
	andata e ritorno che compie il segnale emesso dopo aver rimbalzato sull'oggetto. Sono tipicamente usate per le diagnosi in ambito medico e possono
	essere a due o tre dimensioni.
	\item \textbf{Immagini termiche}: Codificano le proprietà termiche di una determinata scena, definito un intervallo di valori relativo alla
	temperatura. Generalmente usate per la videosorveglianza o nell'edilizia per rilevare eventuali dispersioni di calore.
	\item \textbf{Immagini satellitari}: Categoria generale di immagini bidimensionali che consentono di avere una vista particolarmente ampia su una
	zona specifica della Terra.
	\item \textbf{Immagini radar}: Acquisite tramite elaborazione dei segnali di onde radio; consentono di mostrare oggetti o proprietà di una scena
	che normalmente sarebbero invisibili. Sono infatti usate spesso in ambito meteorologico per indicare valori di umidità o altro.
\end{itemize}

%

\section{Richiami e proprietà dei segnali 1D}
Richiamiamo alcuni concetti notevoli riguardanti i segnali monodimensionali affrontati nel corso dei sistemi. Anzitutto, come già menzionato, un
segnale analogico è una funzione continua che rappresenta la variazione di una quantità fisica in base a una variabile indipendente. Si compone di:
\begin{itemize}
	\item \textbf{Tempo} $t$: Rappresentato nell'asse orizzontale.
	\item \textbf{Ampiezza} $A$: Rappresentata nell'asse verticale e misurata con una quantità predefinita.
	\item \textbf{Frequenza} $f$: Indica la rapidità delle oscillazioni.
\end{itemize}

\noindent La funzione del segnale analogico nel tempo continuo $x(t)$ è definita per ogni valore di tempo $t\in\mathbb{R}$; dopo il campionamento passa
a tempo discreto $x[n]$, quindi la variabile indipendente è definita esclusivamente per valori interi $n\in\mathbb{Z}$. Il segnale digitale è infine
ottenuto con la quantizzazione, la quale discretizza a interi i valori dell'ampiezza all'interno degli intervalli campionati.\par
Seguono ora alcune nozioni e richiami riguardanti i numeri complessi, poiché, applicati ai segnali, sono spesso usati come strumenti matematici per la
semplificazione dell'analisi dei segnali reali, sistemi lineari e sinusoidi. Definiamo quindi un segnale complesso come: \[
    z(t) = x(t) + jy(t)
\]

\noindent dove $x(t) = Re\{z(t)\}$ è la parte reale, $y(t) = Im\{z(t)\}$ quella immaginaria, e $j$ l'unità immaginaria, per la quale vale $j^2 = -1$.
Queste funzioni presentano anche il \textbf{complesso coniugato}, che inverte il valore della parte immaginaria ed è definito come \[
    z(t) = x(t) - jy(t)
\]

\noindent Definiamo inoltre il modulo come la distanza tra l'origine del piano complesso e il punto $P(x(t), y(t))$, calcolata con: \[
    |z(t)| = \sqrt{x^2(t) + y^2(t)}
\]

\noindent La fase è invece l'angolo formato dal vettore complesso con l'asse dei reali, ed è definito geometricamente con la funzione atan2: \[
    argz(t) = \phi(t) = atan2(y(t), x(t))
\]

\noindent Un segnale complesso può essere inoltre rappresentato in forma polare, ovvero: \[
    z(t) = |z(t)|e^{j\phi(t)}
\]

\noindent Le cui coordinate sul piano complesso sono espresse nelle seguenti forme: \[ \begin{cases}
    x(t) = |z(t)|cos(\phi(t))\\
    y(t) = |z(t)|sin(\phi(t))
\end{cases} \]

\noindent In questi termini è possibile esprimere un segnale complesso genericamente come: \begin{equation}
    \begin{split}
        z(t) &= x(t) + jy(t)\\
            &= |z(t)|cos(\phi(t)) + j|z(t)|sin(\phi(t))\\
            &= |z(t)|(cos(\phi(t)) + jsin(\phi(t)))
    \end{split}
\end{equation}

\noindent E, grazie alla \textbf{Formula di Eulero} $e^{j\phi} = cos(\phi) + jsin(\phi)$, otteniamo la scrittura del segnale complesso nella seguente
forma: \[
    z(t) = |z(t)|e^{j\phi(t)}
\]

\noindent A partire da questa formulazione, ci è reso possibile definire una tassonomia dei segnali monodimensionali come segue:

\paragraph{Segnale complesso esponenziale:} $z(t) = e^{jy(t)}$\par
Funzione con un esponente puramente immaginario, con $y(t)$ che ne rappresenta l'omonima componente di un numero complesso generico $z(t)$. Questa
quantità immaginaria può essere interpretata come un angolo: \[
    e^{jy(t)} = cos(y(t)) + jsin(y(t))
\]

\noindent Dalla scrittura otteniamo le parti: \[ \begin{cases}
    Re\{z(t)\} = cos(y(t))\\
    Im\{z(t)\} = sin(y(t))
\end{cases} \]

\noindent E nel piano complesso, la funzione $z(t)$ corrisponde al punto descritto da: \[
    P(t) = (cos(y(t)), sin(y(t)))
\]

\noindent Questo tipo di segnale ha lunghezza unitaria e forma l'angolo $y(t)$ con l'asse dei reali positivo.

\paragraph{Segnali periodici:} $x(t+T) = x(t) \forall t$\par
I segnali periodici a tempo continuo sono tali se esiste un numero positivo $T$, detto \textbf{periodo} e misurato in secondi, tale per cui valga la
definizione. Il segnale si ripete dopo ogni intervallo $T$ e il valore positivo più piccolo che soddisfa questa relazione è detto \textbf{periodo
fondamentale}. Definiamo ora: \begin{itemize}
    \item \textbf{Frequenza} $f$: Numero di cicli completati dal segnale in un secondo, misurata in Hertz. \[
        f = \frac{1}{T}[s]\quad w = 2\pi f = \frac{2\pi}{T}[rad/s]
    \]
    \item \textbf{Ampiezza} $A$: Descrive quanto il segnale devia da uno specifico punto di riferimento, come il suo valor medio. Rappresenta graficamente i picchi della funzione. \[
        x(t) = Acos(\omega t + \phi)\quad A_{pp} = 2|A|
    \]
\end{itemize}

\paragraph{Segnali sinusoidali:} $x(t) = Acos(\omega t + \phi_0)$\par
I suoi parametri sono tutti quelli discussi finora: ampiezza $A$, frequenza angolare $\omega = 2\pi f$, fase iniziale $\phi_0$, fase totale della
sinusoide $\omega t + \phi_0$ e periodo $T$.\par
Il valore di fase iniziale è in particolar modo importante, poiché rappresenta il segnale in un determinato istante del tempo; genericamente è quindi
posto a $\phi(0)$, che equivale a dire $t = 0$. Modificando il suo valore, quindi facendo un cambio di fase, si ottiene un \textbf{time-shift}, una
traslazione temporale rispetto al segnale. Formalmente: \[
    Acos(\omega t + \phi_0) = Acos\left[\omega \left( t + \frac{\phi_0}{\omega}\right)\right] \implies \Delta t = -\frac{\phi_0}{\omega}
\]

\noindent Infine, introduciamo il concetto di \textbf{Energia di un segnale} a tempo continuo, definita generalmente come: \[
    E_x = \int_{-\infty}^{+\infty} |x(t)|^2 dt
\]

\noindent In particolare, per segnali a valori reali, vale $|x(t)|^2 = x^2(t)$, mentre per quelli a valori complessi si usa
$|z(t)|^2 = z(t)\overline{z(t)}$. Inoltre, diciamo che un segnale è \textbf{di wnergia} se vale la relazione: \[
    0 < E_x < +\infty
\]

\noindent Bisogna notare tuttavia che, per i segnali sinusoidali, è possibile estendersi indefinitamente nel tempo poiché continuano a oscillare,
mostrando di conseguenza un valore di energia infinita; per questa ragione viene utilizzato un valore di potenza media quando si lavora con questo tipo
di segnali, o più genericamente quelli periodici.
	\chapter{Signal processing}
	\section{Spazio di Hilbert}
Definiamo lo \textbf{spazio di Hilbert} come uno spazio di vettori $\mathcal{H}$, definito su $\mathbb{R}$ o $\mathbb{C}$, il quale ha un prodotto
interno: \[
    \left<.,. \right> : \mathcal{H}\times \mathcal{H} \to \mathbb{C}
\]

\noindent Il prodotto, supposto che sia lineare nel suo primo argomento, soddisfa alcune proprietà: \begin{itemize}
    \item \textbf{Linearità}: Per $z,z',\omega\in \mathcal{H}$ e $\alpha, \beta\in \mathbb{C}$, vale $\left<\alpha z + \beta z', \omega\right> =
    \alpha\left<z, \omega \right> + \beta\left<z', \omega \right>$
    \item \textbf{Simmetria coniugata}: $\left<z, z' \right> = \overline{\left<z', z \right>}$
    \item \textbf{Positività}: $\left<z, z \right> \geq 0$
    \item \textbf{Definitività positiva}: $\left<z, z \right> = 0 \iff z = 0$
\end{itemize}

\noindent Il prodotto interno induce anche una norma $||z|| = \sqrt{\left< z, z\right>}$, rispetto alla quale lo spazio Hilbertiano è completo. Questa
sua caratteristica è anche ciò che lo distingue da altri spazi semplici di prodotti interni. Inoltre, in quanto completo, ogni sequenza di Cauchy $z_n$
di elementi di $\mathcal{H}$, definita come: \[
    ||z_n - z_m|| \to 0\quad \text{con n, m} \to +\infty
\]

\noindent convergerà, secondo la norma, a un elemento appartenente allo spazio $\mathcal{H}$, il quale, per segnali unidimensionali, è considerato come
$\mathcal{H} = L^2(\mathbb{R};\mathbb{C})$, i cui elementi sono funzioni misurabili che associano un valore dei reali a uno nei complessi e soddisfano:
\[
    \int_{-\infty}^{+\infty} |z(t)|^2 dt < +\infty
\]

\noindent Concludendo che un segnale complesso a energia finita appartiene allo spazio dei segnali $z\in L^2(\mathbb{R};\mathbb{C})$, la cui forma
vettoriale è data da: \[
    (\alpha z + \beta z')(t) = \alpha z(t) + \beta z'(t)
\]

\noindent Dove $z, z'$ sono i vettori dei segnali, mentre $z(t), z'(t)$ sono numeri complessi che indicano il valore assunto dal vettore segnale al
tempo $t$.\par
Nel contesto della teoria dei segnali, il \textbf{prodotto interno} dei segnali può essere interpretato come l'\textbf{energia} del segnale. Questo
perché, a partire dalla forma del prodotto interno standard nello spazio hilbertiano, otteniamo la norma: \[
    \left<z, z' \right> = \int_{-\infty}^{+\infty} z(t)\overline{z'(t)} dt \implies ||z|| = \sqrt{\left<z, z \right>}
\]

\noindent Di conseguenza, se rimuoviamo la radice elevando la scrittura al quadrato, otteniamo: \[
    ||z||^2 = \left<z, z \right> = \int_{-\infty}^{+\infty} z(t)\overline{z(t)} dt
\]

\noindent E siccome vale $z(t)\overline{z(t)} = |z(t)|^2$, possiamo concludere la tesi, dimostrando che il prodotto interno è interpretabile come
l'energia del segnale, la cui simmetria è garantita dal complesso coniugato.\par
Negli spazi di Hilbert, quindi anche negli Euclidei, un vettore può essere scritto come la combinazione dei suoi elementi lungo gli assi cartesiani,
ottenuti attraverso proiezioni: \[
    x = x_1e_1 + x_2e_2 + ... \implies x_k = \left< x, e_k\right>
\]

\noindent Dove l'insieme $e_k$ è una base ortonormale per lo spazio, detta \textbf{base canonica}. Similmente, nello spazio hilbertiano abbiamo: \[
    z(t) = \sum_{k}c_k\phi_k(t)
\]

\noindent Dove il segnale generico $z(t)\in \mathcal{H}$ è definito come combinazione lineare tra una base ortonormale $\phi_k$ e i coefficienti $c_k$.
Diciamo inoltre, che la base è ortonormale se: \[
    \left<\phi_k, \phi_m \right> = \delta_{km}\quad \text{con } \delta_{km} = \begin{cases}
        1, k=m\\
        0, k\neq m
    \end{cases}
\]

\noindent Ogni funzione base ha quindi una norma unitaria $||\phi_k|| = 1$ e i coefficienti sono dati da $c_k = \left<z, \phi_k \right>$.

%

\section{Basi di Fourier}
Introdotto lo spazio di lavoro, possiamo definire la \textbf{base di Fourier} $\{\phi_k(t)\}$, la quale è composta dalle funzioni di base: \[
    \phi_k(t) = \frac{1}{\sqrt{T}}e^{j2\pi kt/T}\quad \text{con }k\in\mathbb{Z}
\]

\noindent Per la quale ogni funzione di base è parametrizzata da $\omega_k = \frac{2\pi k}{T}$. Si può inoltre notare come appartenga alla classe di
funzioni sinusoidali complesse $z_k(t) = Ae^{j(\omega_kt+\phi_0)}$, dove: \begin{itemize}
    \item \textbf{Ampiezza}: $A = \frac{1}{\sqrt{T}}$
    \item \textbf{Frequenza angolare}: $\omega_k = \frac{2\pi k}{T}$
    \item \textbf{Fase iniziale}: $\phi_0 = 0$
    \item \textbf{Parte reale}: $Re\{\phi_k(t)\} = \frac{1}{\sqrt{T}}cos\left(\frac{2\pi kt}{T}\right)$
    \item \textbf{Parte immaginaria}: $Im\{\phi_k(t)\} = \frac{1}{\sqrt{T}}sin\left(\frac{2\pi kt}{T}\right)$
    \item \textbf{Modulo}: $|\phi_k(t)| = \frac{1}{\sqrt{T}}$
    \item \textbf{Arg}: $arg(\phi_k(t)) = \omega_kt = \frac{2\pi kt}{T}$
\end{itemize}

\noindent $T$, come già menzionato, è detto periodo fondamentale, per il quale sono definite anche le frequenze fondamentali che ritornano le frequenze
della base di Fourier: \[
    \omega_k = \frac{2\pi k}{T} = k\omega_0\quad f_k = \frac{k}{T} = kf_0
\]

\noindent Nominiamo, in questi termini $k=1$ la \textbf{frequenza fondamentale}, annotata con $\omega_1 = 1\omega_0 = \omega_0$, e con valori di $k$
maggiori di $1$ diremo che sono la seconda, terza, quarta \textbf{armonica}, e così via. Invece, per quanto riguarda valori negativi, il segnale è
rappresentato da una componente dalla frequenza negativa $\omega_{-k} = -\omega_k$. Dunque possiamo dire che $\phi_k(t)$ è la k-esima funzione di base
complessa armonica.\par
La base è interpretabile geometricamente per ogni valore fissato $t$; il numero complesso dato dalla base è un punto sul cerchio $|z| =
\frac{1}{\sqrt{T}}$, e al variare di $t$, ruota il punto intorno all'origine mantenendo un raggio costante uguale a $\frac{1}{\sqrt{T}}$ con posizione
angolare data da $\frac{2\pi kt}{T}$. In particolare: \begin{itemize}
    \item Con $k>0$ si ha una rotazione antioraria.
    \item Con $k<0$ si ha una rotazione oraria.
    \item Con $k=0$ non si ha alcuna rotazione.
\end{itemize}

\noindent Infine, le funzioni di \textbf{base di Fourier} formano una \textbf{base ortonormale} dello spazio $L^2([0,T], \mathbb{C})$, infatti: \[
    \left<\phi_k, \phi_m \right> = \int_{0}^{T} \phi_k(t)\overline{\phi_m(t)} dt = \delta_{km}
\]

\noindent Il risultato è detto \textbf{delta di Kronecker}, e rispetta le proprietà dell'ortonormalità; mostrando contemporaneamente che le funzioni di
base hanno norma unitaria $||\phi_k|| = 1$ e sono mutuamente ortogonali, dunque l’insieme $\{\phi_k\}$ è una base ortonormale e per ogni
$z(t)\in L^2([0,T], \mathbb{C})$ possiamo scrivere: \[
    z(t) = \sum_{k\in \mathbb{Z}} c_k\phi_k(t)
\]

%

\section{Serie di Fourier}
Dato un segnale periodico $z(t)\in L^2([0,T], \mathbb{C})$, è possibile espandere il segnale nella forma discussa pocanzi: \[
    z(t) = \sum_{k = -\infty}^{+\infty} c_k\phi_k(t)
\]

\noindent Dove i coefficienti sono le proiezioni di $x$ sulle funzioni di base di Fourier: $c_k = \left<x, \phi_k \right>$. Sostituendo la definizione
del prodotto interno otteniamo la \textbf{serie di Fourier}: \[
    c_k = \frac{1}{\sqrt{T}} \int_0^T x(t)e^{-j\omega_k t} dt
\]

\noindent Ogni coefficiente $c_k$ indica quanto la frequenza $\omega_k = \frac{2\pi k}{T}$ è presente all'interno del segnale; in altre parole
rappresenta la proiezione del segnale sulla sinusoide di frequenza $\omega_k$\par
Grazie alla serie di Fourier è possibile fare delle inferenze in merito alle componenti della scrittura. In particolare, se consideriamo la componente
k-esima della serie $z_k(t) = c_k\phi_k(t)$ e scriviamo il coefficiente complesso di Fourier in forma polare, otteniamo: \[
    c_k = |c_k|e^{jarg(c_k)}
\]

\noindent E da questo ricaviamo: \begin{equation}
    \begin{split}
        z_k(t) &= c_k\phi_k(t)\\
            &= |c_k|e^{jarg(c_k)}\cdot \frac{1}{\sqrt{T}}e^{j\omega_k t}\\
            &= \frac{|c_k|}{\sqrt{T}}e^{j(\omega_kt + arg(c_k))}
    \end{split}
\end{equation}

\noindent Dunque $z_k(t)$ è una sinusoide complessa della forma $A_ke^{j(\omega_kt + \phi_{0,k})}$, dove: \begin{itemize}
    \item \textbf{Ampiezza}: $A_k = |z_k(t)| = |c_k||\phi_k(t)| = \frac{|c_k|}{\sqrt{T}}$
    \item \textbf{Fase}: $arg(z_k(t)) = \omega_kt + arg(c_k)$
    \item \textbf{Fase iniziale}: $\phi_{0,k} = arg(c_k)$
\end{itemize}

\noindent Inoltre, sapendo che $c_k$ è un numero complesso; $|c_k|$ controlla la scala del componente e $\omega_kt$ fornisce la fase dipendente dal
tempo. In questi termini, la funzione $c_k\phi_k(t)$ è detta componente armonica k-esima del segnale.

%

\section{Trasformata di Fourier}
Quando parliamo di segnali non periodici, il periodo fondamentale $T$ tende idealmente a infinito, avvicinando sempre di più le frequenze discrete
$\omega_k = \frac{2\pi k}{T}$, eventualmente facendo loro formare un insieme continuo. In questi termini, la sommatoria della serie di Fourier diventa
un integrale della forma: \[
    x(t) = \frac{1}{\sqrt{2\pi}}\int_{-\infty}^{+\infty} X(\omega)e^{i\omega t}d\omega
\]

\noindent Dal quale possiamo definire la \textbf{trasformata di Fourier}: \[
    X(\omega) = \frac{1}{\sqrt{2\pi}}\int_{-\infty}^{+\infty} x(t)e^{-i\omega t} dt
\]

\noindent In particolare, la funzione $\phi_\omega(t) = \frac{1}{\sqrt{2\pi}}e^{i\omega t}$ rappresenta le funzioni di base associate alla frequenza
$\omega$. Formalmente si scrive: \[
    X(\omega) = \left<x, \phi_\omega \right>
\]

\noindent E indica che la $X(\omega)$ è la proiezione del segnale sulla sinusoide complessa di frequenza $\omega$. Attenzione, tuttavia, che le funzioni
$e^{i\omega t}$ non appartengono allo spazio $L^2(\mathbb{R}, \mathbb{C})$, dunque, nel caso non periodico, non sono funzioni di base ordinarie dello
spazio di Hilbert, bensì quello che fanno è formare una base generalizzata; per questa ragione è più corretto dire che la trasformata di Fourier è la
\textbf{rappresentazione} di un segnale rispetto a una famiglia continua di funzioni di base generalizzata, come quella delle sinusoidi complesse.\par
In ogni caso, facendo uso della trasformata, è possibile passare dal dominio del tempo, i cui segnali sono in forma $x(t)$ a quello delle frequenze
$X(\omega)$, dove ogni valore è una proiezione: $X(\omega) = \left<x, \phi_\omega \right>$. Per ricostruire il segnale è necessario possedere i valori
di ampiezza $|X(\omega)|$ e fase $arg(X(\omega))$. Inoltre, siccome la trasformata di Fourier è un operatore unitario, è preservato il prodotto interno
e, di conseguenza, la norma: \[
    \int_{-\infty}^{+\infty} |x(t)|^2 dt = \int_{-\infty}^{+\infty} |X(\omega)|^2 d\omega
\]

\noindent Questa scrittura è detta \textbf{relazione di Parseval} sotto convenzione simmetrica, e la quantità $|X(\omega)|^2$ è interpretabile come
densità di energia rispetto alla frequenza, rendendo calcolabile l'energia totale del segnale anche nel dominio delle frequenze. Se si vogliono
considerare solo le frequenze contenute in una certa \textbf{banda} $B$, è possibile moltiplicare lo spettro per una funzione indicatore: \[
    X_B(\omega) = X(\omega)l_B(\omega)
\]

\noindent Che fornirà il segnale filtrato: \[
    x_B(t) = \frac{1}{\sqrt{2\pi}}\int_B X(\omega)e^{i\omega t}d\omega
\]

\noindent Il quale è una proiezione sul subspazio generato dalle sinusoidi le cui frequenze si trovano all'interno della banda. Un filtro ideale è
considerabile come un operatore di proiezione sullo spazio delle frequenze.
	\chapter{Formazione delle immagini}
	


\begin{comment}
	Filtraggio analogico e digitale (filtri passa basso, passa alto, lineari e non lineari) nel dominio del tempo e della frequenza
\end{comment}
	\chapter{Calibrazione}
	

\begin{comment}
	Conversione A/D (campionamento e quantizzazione di segnali e immagini)
\end{comment}
	\chapter{Geometria proiettiva}
	

\begin{comment}
	Segmentazione di immagini (principali algoritmi di estrazione di contorni e analisi basata su regioni)
\end{comment}
	\chapter{Filtri, edge e contorni}
	

\begin{comment}
	Analisi delle tessiture mediante filtri di Gabor e cenni di analisi tempo/frequenza e multiscala
\end{comment}
	\chapter{Segmentazione}
	

\begin{comment}
	Estrazione di feature ed esempi di applicazione a problemi di pattern recognition
\end{comment}
	\chapter{Operatori morfologici}
	\section{Riconoscimento e collegamento AI}

%

\section{Acquisizione e immagini range}

%

\section{Feature detection and matching}
	\chapter{Utilizzo dell'AI}
	\section{Riconoscimento e collegamento AI}

%

\section{Acquisizione e immagini range}

%

\section{Feature detection and matching}

\end{document}